\documentclass[11pt]{article}
\usepackage[margin=1in]{geometry}
\usepackage{amsmath,amssymb,amsthm,booktabs,array}
\usepackage[T1]{fontenc}
\usepackage{lmodern}
\usepackage{microtype}
\usepackage{hyperref}
\hypersetup{colorlinks=true,linkcolor=blue,citecolor=blue,urlcolor=blue}
\setlength{\emergencystretch}{3em}
\newtheorem{theorem}{Theorem}
\newtheorem{proposition}{Proposition}
\newcommand{\norm}[1]{\left\lVert #1\right\rVert}
\newcommand{\cI}{\mathcal I}
\title{A Finite Cutoff-Sensitivity Obstruction for the V59 Residual Phase}
\author{Liang Wang\\School of Mathematics and Statistics\\Huazhong University of Science and Technology (HUST)\\Wuhan, China}
\date{26 August 2026}
\begin{document}
\maketitle
\begin{abstract}
TPC-267 found a finite quarter contraction for a literal V59 residual using
the coarse local comparison cutoff \(z=2\). We perform the next hostile test
while keeping the prime shell, unit masks, deleted diagonal, source
coefficient, kernel operator, and rank-three projection fixed. Outward
rational intervals classify sixteen finite rows. Six matched \(z=2\)
controls contract below \(1/4\), while six perturbed rows are certified
above \(1/4\). The same central clock
\((N,H,Q,s)=(64,15,4,1)\) changes from
\(\rho\leq0.2320126753\) at \(z=2\) to
\(\rho\geq0.2735126949\) at \(z=3\). This refutes a universal quarter claim
over the declared finite parameter family, but it is not an asymptotic
counterexample and pays no fixed-power credit.
\end{abstract}
\section{Position and claim firewall}
The typed end-to-end audit of TPC-266 requires an actual growing V59
radius or signed-phase estimate before the endpoint budget can be paid.
TPC-267 supplied a reproducible finite interface and twelve \(z=2\) rows
with \(\lvert C_\perp\rvert/R<1/4\), while explicitly leaving the growing
radius and source-specified profile open. The present paper asks whether that
finite inequality is stable under the interface choices declared as modeling
choices.
\begin{center}
\small\texttt{CUTOFF\_SENSITIVITY\_OBSTRUCTION}\\
\textit{numerically certified finite result}.
\end{center}
\section{The finite physical operator}
Let \(N\) be even and set
\[
 \cI_N=\{N/2+1,\ldots,N\},\qquad
 \mathcal Q_Q=\{q:q\ {\rm prime},\ Q<q\leq2Q\}.
\]
For an integer \(z\geq2\), define
\[
 b_N^{(z)}(u)=C_2^{(z)}\,1_{2\nmid u}
 \prod_{p\leq z}\frac p{p-1}
 \prod_{\substack{p\mid u\\p>z}}\frac{p-1}{p-2},       \tag{1}
\]
with value zero when \(u+2\) has a prime factor at most \(z\), and
\[
 C_2^{(z)}=\prod_{p>z}\left(1-\frac1{(p-1)^2}\right).
\]
The source coefficient and residual input are
\[
 \beta_N(t)=\frac{\Lambda(t)}{\log t}
 -\sum_{\substack{d\mid t\\d^{400}\leq N^{133}}}\mu(d),
 \qquad w_N(u)=\Lambda(u+2)-b_N^{(z)}(u).                 \tag{2}
\]
For \(s\in\{1,2,3\}\), use
\[
 K_{H,s}(h)=\left(1+(h/H)^2\right)^{-s},                 \tag{3}
\]
and the literal centered operator
\[
 A(u,t)=1_{u\ne t}\sum_{q\in\mathcal Q_Q}qK_{H,s}(u-t)
 1_{q\nmid ut}\left(1_{u\equiv t\pmod q}-\frac1{q-1}\right). \tag{4}
\]
The prime-only shell, both unit masks, and the deleted diagonal are therefore
part of every row. Put \(g=A\beta_N\).

\section{Projection and interval decision}
Partition \(\cI_N\) into four consecutive blocks of size \(B=N/8\).
Use the three orthogonal contrasts
\[
 c_0=(1,1,-1,-1),\quad c_1=(1,-1,0,0),\quad
 c_2=(0,0,1,-1).                                          \tag{5}
\]
For block sums \(W_j,G_j\), define
\[
 C=\langle w_N,g\rangle,\quad
 C_3=\frac{W_0G_0}{4B}+\frac{W_1G_1}{2B}+\frac{W_2G_2}{2B},
 \quad C_\perp=C-C_3,                                     \tag{6}
\]
and
\[
 R^2=\norm{(I-P_3)w_N}^2\norm{(I-P_3)g}^2,\qquad
 \rho^2=\frac{\lvert C_\perp\rvert^2}{R^2}.                \tag{7}
\]

\begin{proposition}[finite identity]
For every listed row, (6) is the exact cross-Gram decomposition associated
with the rank-three block-contrast projection, and \(R^2>0\).
\end{proposition}

\begin{proof}
The three vectors in (5) have squared norms \(4B,2B,2B\) and are pairwise
orthogonal. The orthogonal projection formula gives (6), and subtracting
the projected cross-Gram gives \(C_\perp\). The two residual norm factors
are obtained by applying the same formula to each squared norm. Positivity
is enclosed by the interval computation described below.
\end{proof}

The finite Euler product is multiplied through \(P=50000\). The tail uses
\[
 \prod_{p>P}\left(1-\frac1{(p-1)^2}\right)
 \geq 1-\sum_{p>P}\frac1{(p-1)^2}
 >1-\frac1{P-1}.                                         \tag{8}
\]
The logarithms are evaluated at 100 decimal digits with a \(10^{-25}\)
guard. Every endpoint is then rounded outward to a rational grid of width
\(10^{-30}\). Thus the decision in (7) uses only rational interval
operations:
\[
 \rho^2_{\rm hi}<\frac1{16}\Rightarrow\text{contraction},\qquad
 \rho^2_{\rm lo}>\frac1{16}\Rightarrow\text{obstruction}. \tag{9}
\]
\section{The finite obstruction}

\begin{theorem}[cutoff sensitivity]
Among the sixteen rows in the certificate, ten are certified contractions
and six are certified obstructions. In particular,
\[
\begin{array}{c|c|c}
(N,H,Q,s,z)&\rho^2\text{ interval}&\text{classification}\\ \hline
(64,15,4,1,2)&[0.0538214595269,\,0.0538298814898]&\text{contraction}\\
(64,15,4,1,3)&[0.0748091943191,\,0.0748218869170]&\text{obstruction}
\end{array}
\]
so changing only the finite cutoff from \(z=2\) to \(z=3\) flips the
quarter classification. The \(z=3\) obstruction persists at \(H=13\)
and \(H=17\), and the row \((64,15,4,1,5)\) has
\(\rho\leq0.3851247936\).
\end{theorem}

\begin{proof}
The producer enumerates (1)--(4), computes the exact projection split, and
propagates the outward intervals through (7). The lower endpoints of the
six obstruction rows are all larger than \(1/16\), while the upper endpoints
of the ten contraction rows are all smaller than \(1/16\). The independent
replay recomputes the same sixteen ratios without importing the producer.
The central pair and the clock-neighborhood rows are therefore separated
from the threshold by certified intervals, not by rounded display values.
\end{proof}

\begin{table}[t]
\centering
\caption{Selected rows. Displayed \(\rho\) values are rounded square-root
upper endpoints from the certificate; the contraction/obstruction decision
uses the separately stored interval for \(\rho^2\).}
\small
\begin{tabular}{rrrrrll}
\toprule
N&H&Q&s&z&displayed \(\rho\)&class\\
\midrule
64&15&4&1&2&0.2320126753&contract\\
96&20&5&1&2&0.2138888503&contract\\
128&24&5&1&2&0.1267919366&contract\\
192&32&6&1&2&0.0066519390&contract\\
256&38&6&1&2&0.0093009287&contract\\
384&50&7&1&2&0.0631257787&contract\\
64&13&4&1&3&0.2689314328&obstruct\\
64&15&4&1&3&0.2735358969&obstruct\\
64&17&4&1&3&0.2758624715&obstruct\\
96&20&5&1&3&0.2996937679&obstruct\\
96&20&5&2&3&0.2886875882&obstruct\\
64&15&4&1&5&0.3851247936&obstruct\\
\bottomrule
\end{tabular}
\end{table}

\section{Interpretation}

The central pair gives a clean finite sensitivity statement. The physical
matrix, source coefficient, clock, shell, masks, and projection are
identical in the two rows; only the declared local comparison cutoff changes.
The interval for the first row lies below \(1/16\), whereas the lower
endpoint for the second row is

\[
 0.0748091943191>1/16=0.0625.
\]

Consequently a universal quarter-sector lemma over this tested finite
parameter family is false. The clock perturbations \(H=13\) and \(H=17\)
show that the \(z=3\) obstruction is not a single rounded-clock accident,
and the \(z=5\) row supplies a stronger finite stress point. These are
finite statements about the explicitly enumerated rows.

There are two important qualifications. First, \(z\) is a modeling
parameter here. The source-level V59 construction has a growing local
cutoff, and this paper does not prove that its finite values are represented
by any fixed \(z\) in the table. Thus the result is not a source-level
growing-\(x\) counterexample. Second, the kernel exponent is not monotone in
this diagnostic: the \(N=64,z=3\) row with \(s=2\) is just below the quarter
threshold, while the \(N=96,z=3\) row with \(s=2\) is above it. A finite
contraction therefore cannot be promoted by an informal smoothness or
monotonicity argument.

The audit also does not pay any endpoint budget. It proves no bound for the
size of \(R\), gives no fixed-power saving, and supplies no arithmetic
\(L^2\) estimate or full Gate-B closure. In the notation of the upstream
firewall, the fixed-power credit remains zero and the twin-prime conclusion
remains open.

\section{Conclusion}

TPC-268 turns the finite residual census into an adversarial diagnostic. A
matched \(z=2\) control and \(z=3\) perturbation use the same finite physical
operator but land on opposite sides of the quarter threshold, with six
outward-certified obstructions across the sixteen-row family. The natural
next question is not whether the most favorable finite profile can be made
to contract, but whether a growing cutoff and a source-compatible smooth
profile admit a uniform interval theorem. Until that question is answered,
the observed phase contraction is evidence about a finite interface rather
than progress toward the asymptotic twin-prime estimate.

\begin{thebibliography}{9}
\bibitem{hardywright}
G.~H.~Hardy and E.~M.~Wright,
\emph{An Introduction to the Theory of Numbers}, 6th ed., Oxford
University Press, 2008.

\bibitem{montgomeryvaughan}
H.~L.~Montgomery and R.~C.~Vaughan,
\emph{Multiplicative Number Theory I: Classical Theory}, Cambridge
University Press, 2007.

\bibitem{tpc267}
L.~Wang, ``A finite literal V59 residual-radius and signed-phase census,''
TPC-267 project release, 2026.
\end{thebibliography}

\end{document}
